\documentclass[twoside]{article}
\usepackage{package}
\usepackage[T1]{fontenc}
\usepackage{subfiles}
% Silenciar la advertencia de versalita + negrita inexistente
\WarningFilter{latex}{Some font shapes were not available}
\WarningFilter{latex}{Font shape `T1/lmr/bx/sc'}

% Silenciar la advertencia de caracteres perdidos en el nullfont (para 01_introduccion.tex)
\WarningFilter{latex}{Missing character: There is no}
%Nuevos comandos
\renewcommand{\thefootnote}{\arabic{footnote}}
\raggedbottom



\begin{document}
\begin{center}
    \thispagestyle{empty}
    \large{\textsc{Integración en MATLAB y Sistemas Lineales}}\\[0.5cm] 
    \large{Fabrizio Benavidez\footnote[1]{\texttt{fabrizio.benavidez@mi.unc.edu.ar}}}\\[0.5cm]

     \href{https://github.com/Fabopy/Matlab-curso-}{%
        {\Large\faGithub}\quad
        \nolinkurl{https://github.com/Fabopy/Matlab-curso-}%
    }\\[0.5cm]

    \large{\footnotemark[1]\textit{Facultad de Ciencias Exactas, Físicas y Naturales, Ingeniería Aeronáutica}}
    \\[0.5cm]

    \begin{minipage}{0.7\textwidth}
    \hrule
    \vspace{0.2cm}
        \normalsize{\textsc{Resumen} En el siguiente trabajo se realiza integración numérica en MATLAB mediante algunos métodos propuestos. Se utiliza la regla del Trapecio, Simpson y además la función propia de MATLAB para integración. Además, se plantea el problema de resolver la integración como ecuación diferencial de primer grado, utilizando funciones de Runge-Kutta para su solución. 
        
        Por otro lado, se estudia el caso de un brazo robótico para estimar su función de transferencia, proponiendo un sistema de segundo orden.}\\[0.2cm]
        \normalsize{\textit{Palabras claves}: Integración, MATLAB, Modelo\ }
    \vspace{0.2cm}
    \hrule
    \end{minipage}
    
    \vfil

    \includegraphics[width = 0.4\textwidth]{imagenes/Matlab_Logo.png}\\[0.5cm]

    \vfil

    \today\\[0.2cm]
    \large{Semana 11}
\end{center}







\newpage

\section*{Guía 21: Integración en MATLAB}
\subfile{secciones/guia21}


\newpage

\section*{Guía 22: Identificadores de Sistemas Lineales} 
\subfile{secciones/guia22}














\newpage
% \thispagestyle{empty}
% \bibliography{biblioteca.bib}
% \bibliographystyle{plain}


\end{document}